% =====================================================================
% cap-01-fundamentacao.tex — Capítulo 2: Fundamentação Teórica
% =====================================================================

\chapter{Fundamentação Teórica}
\label{cap:fundamentacao}

Este capítulo apresenta os conceitos teóricos fundamentais que sustentam
o sistema desenvolvido, organizados em quatro blocos: o sensor de
profundidade Kinect (Seção~\ref{sec:kinect}), o modelo de câmera
\textit{pinhole} estendido por Tsai (Seção~\ref{sec:pinhole}), as
ferramentas de álgebra linear computacional empregadas
(Seção~\ref{sec:algebra}) e a representação cartográfica por Modelos
Digitais de Elevação (Seção~\ref{sec:mde}).

\section{O Sensor Microsoft Kinect}
\label{sec:kinect}

O Microsoft Kinect é um dispositivo de captura tridimensional
originalmente concebido para o controle gestual do videogame
Xbox 360 (versão~v1, 2010) e posteriormente sucedido pela
versão para Xbox One (versão~v2, 2014). Para fins acadêmicos e
industriais, o sensor revelou-se notável pela combinação de baixo custo
e qualidade aceitável de profundidade, tendo se consolidado como
plataforma de pesquisa em visão computacional, robótica e interação
humano-computador~\cite{zhang2012kinect}.

\subsection{Princípio de medição}
\label{subsec:tof}

A versão~v1 do Kinect mede a profundidade pelo princípio de
luz estruturada: um projetor infravermelho emite um padrão
pseudo-aleatório de pontos sobre a cena, e uma câmera infravermelha
calibrada captura a deformação desse padrão. A profundidade de cada
\emph{pixel} é estimada por triangulação. A versão~v2 substitui esse
princípio pelo \textit{Time-of-Flight} (ToF), no qual a profundidade é
calculada a partir do deslocamento de fase de pulsos infravermelhos
modulados~\cite{sell2014xbox}.

\subsection{Parâmetros intrínsecos}
\label{subsec:intrinsecos-kinect}

Para a versão~v2, empregada como referência neste trabalho, os
parâmetros intrínsecos do sensor de profundidade são:

\begin{table}[!ht]
\centering
\caption{Parâmetros intrínsecos do sensor de profundidade do Kinect~v2.}
\label{tab:intrinsecos}
\begin{tabular}{ll}
\toprule
\textbf{Parâmetro} & \textbf{Valor} \\
\midrule
Resolução                  & $512 \times 424$ pixels \\
Distância focal $f_x, f_y$ & $367{,}35$ pixels \\
Centro óptico $(c_x, c_y)$ & $(260{,}00,\; 205{,}00)$ pixels \\
Faixa útil de operação     & $0{,}5$ a $4{,}5$ metros \\
Saída do SDK               & milímetros (\texttt{uint16}) \\
\bottomrule
\end{tabular}
\end{table}

A conversão entre pixel $(u, v)$ e ponto 3D $(X, Y, Z)$ no referencial
do sensor usa a relação canônica de \textit{back-projection} pinhole:

\begin{equation}
\label{eq:backproj}
X = (u - c_x) \cdot \frac{Z}{f_x}, \qquad
Y = (v - c_y) \cdot \frac{Z}{f_y}, \qquad
Z = \frac{d_{\text{mm}}}{1000},
\end{equation}

\noindent onde $d_{\text{mm}}$ é a profundidade lida pelo sensor.

\subsection{Sensores alternativos: Azure Kinect e Intel RealSense}
\label{subsec:sensores-alt}

O Microsoft Kinect~v2 foi descontinuado pela fabricante em 2015, o que
motiva a avaliação de sensores substitutos para a continuidade da linha
de pesquisa. O Azure Kinect DK, sucessor direto lançado em
2019, mantém o princípio \textit{Time-of-Flight} da versão~v2, porém com
maior resolução de profundidade ($1024 \times 1024$ no modo
\textit{narrow field of view}) e uma unidade de medição inercial (IMU)
integrada, útil para compensar pequenas oscilações do suporte do sensor.
A família Intel RealSense (modelos D400), por sua vez, emprega
visão estéreo ativa: um par de câmeras infravermelhas associado
a um projetor de padrão estrutural estima a profundidade por
triangulação, dispensando o cálculo de fase do \textit{Time-of-Flight} e
reduzindo o consumo energético, ao custo de maior ruído em superfícies
pouco texturizadas — não é o caso da areia, cuja textura granular
favorece a correspondência estéreo.

Ambos os sensores expõem SDKs distintos do PyKinect2 utilizado para o
Kinect~v2, mas convergem para o formato de nuvem de pontos RGB-D
consumido pela biblioteca \texttt{Open3D}~\cite{zhou2018open3d}, já
integrada à cadeia de \textit{fallback} do sistema
(Seção~\ref{subsec:fallback}) como segundo nível, imediatamente abaixo
do PyKinect2. Na prática, isso significa que a migração do Kinect~v2
para um desses sensores mais recentes exige apenas a captura da nuvem no
formato esperado por \texttt{Open3D} — todo o restante do
\textit{pipeline} matemático (Capítulo~\ref{cap:motor}), agnóstico ao
fabricante do sensor, permanece inalterado. Essa possibilidade é
retomada como trabalho futuro na Conclusão.

\subsection{Ruído e tratamento}
\label{subsec:ruido}

O sensor apresenta ruído gaussiano cujo desvio padrão cresce
quadraticamente com a profundidade. Na faixa de operação adotada
(2--3~m, com o sensor a 2,5~m da mesa), valores típicos situam-se entre
$\sigma \approx 2$ e $\sigma \approx 4$~mm por \textit{pixel}. Esse
ruído justifica a estratégia de agregação espacial
(Capítulo~\ref{cap:rendering}): ao se calcular a média de $n$ pontos por
célula da malha, o desvio padrão do estimador da altura é reduzido
para $\sigma_{\bar{Z}} = \sigma / \sqrt{n}$, conforme o
\textit{Central Limit Theorem}.

\section{Modelo de Câmera Pinhole e Calibração de Tsai}
\label{sec:pinhole}

A projeção de um ponto 3D no plano de imagem 2D segue, em primeira
aproximação, o modelo \textit{pinhole}:

\begin{equation}
\label{eq:pinhole}
s
\begin{bmatrix} u \\ v \\ 1 \end{bmatrix}
=
K \cdot
\begin{bmatrix} R_{\text{ext}} & \mathbf{t}_{\text{ext}} \end{bmatrix}
\begin{bmatrix} X \\ Y \\ Z \\ 1 \end{bmatrix},
\quad
K = \begin{bmatrix} f_x & 0 & c_x \\ 0 & f_y & c_y \\ 0 & 0 & 1 \end{bmatrix},
\end{equation}

\noindent onde $K$ é a matriz intrínseca do dispositivo (câmera ou
projetor), $R_{\text{ext}}$ e $\mathbf{t}_{\text{ext}}$ formam a
matriz extrínseca (rotação e translação do referencial mundo
para o referencial do dispositivo) e $s$ é um fator de escala não nulo.

O modelo de Tsai~\citeonline{tsai1987versatile} estende o
\textit{pinhole} canônico com a inclusão de distorções óticas
radiais (coeficientes $k_1, k_2, k_3$) e tangenciais ($p_1, p_2$),
fornecendo precisão sub-pixel para câmeras e projetores reais. A
biblioteca OpenCV~\cite{bradski2008opencv} implementa esse modelo
através das funções \texttt{cv2.calibrateCamera}, \texttt{cv2.solvePnP}
e \texttt{cv2.projectPoints}, todas utilizadas no presente trabalho.

\section{Álgebra Linear Computacional}
\label{sec:algebra}

\subsection{Decomposição em Valores Singulares}
\label{subsec:svd}

Dada uma matriz $M \in \mathbb{R}^{N \times 3}$, sua decomposição em
valores singulares é $M = U \Sigma V^{T}$, com $\Sigma$ diagonal contendo
$\sigma_1 \geq \sigma_2 \geq \sigma_3 \geq 0$. O vetor singular direito
associado ao menor valor $\sigma_3$, isto é, a última linha de $V^{T}$,
constitui a solução do problema de mínimos quadrados restrito
\begin{equation}
\label{eq:svd-min}
\mathbf{n}^* = \arg\min_{\|\mathbf{v}\| = 1} \; \|M \mathbf{v}\|_2^2,
\end{equation}

\noindent o que, no contexto deste trabalho, corresponde ao vetor normal
do plano que melhor se ajusta à nuvem de pontos centralizada
(Seção~\ref{sec:passo1}).

\subsection{Amostragem Robusta: RANSAC}
\label{subsec:ransac}

O ajuste de plano por SVD (Seção~\ref{subsec:svd}) é um estimador de
mínimos quadrados: minimiza a soma dos quadrados das distâncias
de \emph{todos} os pontos da nuvem ao plano, e por isso é sensível a
\textit{outliers} — pontos que não pertencem à superfície de interesse.
No caso de campo deste trabalho, o campo de visão do Kinect é mais largo
que a área útil do caixão de areia, de modo que a nuvem capturada durante
a calibração inclui, além da superfície de referência, a moldura de
madeira do caixão, o piso ao redor e ruído da sala — outliers que, se
incluídos ingenuamente no SVD, distorcem a normal calculada.

O algoritmo RANSAC (\textit{Random Sample Consensus}), proposto
por Fischler e Bolles~\citeonline{fischler1981random}, resolve esse problema por
amostragem aleatória repetida. Dada uma nuvem de $N$ pontos, cada
iteração executa três passos: (i) sorteiam-se $k$ pontos mínimos
necessários para definir o modelo — no caso de um plano, $k = 3$; (ii)
constrói-se o modelo candidato a partir desses pontos; (iii) conta-se o
número de inliers — pontos da nuvem inteira cuja distância ao
modelo candidato é inferior a um limiar $d_{\text{in}}$. Após um número
fixo de iterações, o modelo com maior número de inliers é escolhido como
vencedor. Diferentemente do SVD, o RANSAC não assume que a maioria dos
pontos seja bem descrita por um único modelo suave; ele explicitamente
\emph{separa} a nuvem em inliers e outliers, sendo tolerante a frações
significativas de dados espúrios, desde que o conjunto de inliers
permaneça majoritário dentro da amostra observada.

Neste trabalho, o RANSAC não substitui o SVD, mas o precede: o
plano vencedor do RANSAC delimita o conjunto de inliers (presumivelmente
a superfície de referência plana), e apenas esses pontos alimentam o
refinamento por mínimos quadrados via SVD, que produz a normal final com
precisão sub-milimétrica. Essa combinação — robustez do RANSAC seguida
da precisão do SVD sobre os inliers selecionados — é detalhada, com a
formalização algébrica completa e sua aplicação à calibração
Kinect~$\to$~Mesa, na Seção~\ref{sec:passo1}.

\subsection{Processo de Gram-Schmidt}
\label{subsec:gs}

O processo de ortogonalização de Gram-Schmidt converte um conjunto
linearmente independente $\{\mathbf{v}_1, \ldots, \mathbf{v}_n\}$ em um
conjunto ortonormal preservando os subespaços gerados. Neste trabalho,
o procedimento é aplicado para construir uma base ortonormal
$\{X_{\text{mesa}}, Y_{\text{mesa}}, Z_{\text{mesa}}\}$ a partir do
vetor normal obtido por SVD.

\subsection{Transformações afins homogêneas}
\label{subsec:afim}

Uma transformação afim 3D é representada de forma compacta por uma
matriz $4 \times 4$ que combina rotação $R \in \mathbb{R}^{3 \times 3}$ e
translação $\mathbf{t} \in \mathbb{R}^{3}$:
\begin{equation}
\label{eq:afim}
T = \begin{bmatrix} R & \mathbf{t} \\ \mathbf{0}^{T} & 1 \end{bmatrix},
\qquad
\begin{bmatrix} \mathbf{p}' \\ 1 \end{bmatrix}
= T \begin{bmatrix} \mathbf{p} \\ 1 \end{bmatrix}.
\end{equation}

A composição de transformações é dada pelo produto matricial. Em
particular, a aplicação sucessiva de $T$ (alinhamento) e $T_{\text{shift}}$
(deslocamento ao centro lógico da mesa) é realizada por
$T_{\text{final}} = T_{\text{shift}} \cdot T$, conforme
discutido na Seção~\ref{subsec:T-shift}.

\section{Modelos Digitais de Elevação e o Formato GeoTIFF}
\label{sec:mde}

Um Modelo Digital de Elevação (MDE) é uma representação
matricial discreta da altimetria de uma região, na qual cada \emph{pixel}
codifica a elevação $Z$ correspondente a uma coordenada geográfica
$(X, Y)$~\cite{li2005digital}. MDEs constituem produto de saída
fundamental da Geomática e são amplamente empregados em planejamento
militar, hidrografia, agricultura de precisão e simulação ambiental.

O formato GeoTIFF~\cite{geotiff} é um padrão aberto que
estende o tradicional \texttt{TIFF} com metadados de georreferenciamento:
sistema de coordenadas, \emph{datum}, projeção cartográfica e
transformação afim que mapeia índices de \emph{pixel} para coordenadas
do mundo. A leitura programática do formato em Python é realizada pela
biblioteca \texttt{rasterio}~\cite{rasterio}, que retorna um arranjo
NumPy associado a metadados consultáveis. A consulta da elevação em
coordenadas arbitrárias entre os pixels é tratada neste trabalho por
interpolação bilinear, implementada na classe
\texttt{scipy.interpolate.RegularGridInterpolator}~\cite{scipy}.
